\clearpage
\fancyhead[L]{\small EINLEITUNG}
\SourceLeaf{351}
\begin{center}
{\Large EINLEITUNG}\par\medskip
{\small [FÜR DIE ZEITSCHRIFT: RESULTATE AUS DEN BEOBACHTUNGEN DES MAGNETISCHEN VEREINS IM JAHRE 1836. HERAUSGEGEBEN VON C. F. GAUSS UND W. WEBER.]}
\end{center}
\medskip
Unter den mannigfaltigen Äusserungen der erdmagnetischen Kraft, deren Ergründung nur durch zahlreiche an den verschiedensten Punkten der Erdoberfläche fortgesetzt anzustellende genaue Beobachtungen zu erreichen ist, bedürfen die \textit{unregelmässigen Änderungen}, welchen wir jene Kraft unterworfen finden, am meisten eines streng geordneten Zusammenwirkens der Beobachter. Es ist bekannt genug, dass die Bestimmungsstücke der erdmagnetischen Kraft, die Abweichung, die Neigung und ohne Zweifel auch die Stärke (wenn gleich in Beziehung auf die letzte, die erst seit einigen Jahrzehnden in den Kreis der Forschungen aufgenommen ist, noch hinlängliche Erfahrungen fehlen) fortwährend Veränderungen erleiden, seculäre erst nach längerer Zwischenzeit in die Augen fallende, aber im Laufe der Zeit sehr beträchtlich werdende, und periodische nach den Jahres- und Tageszeiten wechselnde. Aber so weit diese Veränderungen mit Regelmässigkeit geschehen, ist ein streng geordnetes Zusammenwirken der Beobachter an verschiedenen Orten, wenn auch für die Beschleunigung der Erweiterung unserer Einsicht höchst wünschenswerth, doch nicht wesentlich nothwendig, und jeder Beobachter kann auch unabhängig von den andern nützliche Beiträge liefern.

Anders verhält es sich hingegen mit den unregelmässigen Veränderungen, welchen man erst in den letzten Jahren eine grössere Aufmerksamkeit zu widmen angefangen hat. Dass während der Sichtbarkeit eines Nordlichts die Magnetnadel unregelmässige und oft sehr grosse Bewegungen zeigt, haben schon vor bei%
